\section{多场战争挤在同一条粮道上}

1851年金田起事后，太平军北上，1853年占领南京并建都天京。长江城市、粮道和税源随之成为战场；北伐、西征、上海小刀会、捻军扩展又把华北、江淮和租界城市带进不同的战争节奏。1856年天京事变削弱太平天国的协调能力，同年第二次鸦片战争开始，清廷不得不在内战与海防之间分配本已不足的兵饷和注意力。1858年条约、1860年北京危机不是内战的插曲，而是重新规定财政、港口和外交机构的压力来源。

危机的关键不只是“中央弱、地方强”。湘军、淮军、团练与地方筹饷是在训练、家族网络、绅商捐输、厘金和军需转运中形成的；它们提高了清廷在若干战区的作战能力，也让军队、税收和任命更深地嵌入省级关系。1864年天京陷落终结了太平天国的中心政权，却没有使难民、荒地、债务和地方驻军同时消失。1868年捻军主力被击败后，战争留下的筹饷制度仍在塑造国家。

1855年黄河改道尤能显示战争并非纯粹的军事故事。河道、堤防、迁徙、粮价与军队移动相互影响；铜瓦厢的一次决口既不能解释全部战局，也不能被放在战事之外。地方家庭要在征发、逃难、守土、借贷和重建之间作出有限选择，正是这些选择使省份成为战后政治的基础。

\subsection{材料与争议}

军机档、方略和实录可追踪任命、城池和清廷的军费命令；太平文献、地方志、契约、日记、租界报刊和传教士材料各有立场与地域限制。兵数、伤亡、地方支持和战争损失不宜接受任何一方战报的总计，应将军事行动与社会后果分层叙述。
